% Supplementary details for Section~\ref{sec:synthetic_data}.

\paragraph{Global settings and splits.}
All simulations start with $G=1000$ genes and are filtered to retain genes
expressed in at least three observed cells. We generate three splits:
(\textit{i}) a tuning split with $N=1000$ cells per scenario,
(\textit{ii}) a held-out test split with $N=5000$ cells per scenario, and
(\textit{iii}) a runtime-scaling split with
$N\in\{10000,20000,\ldots,100000\}$ on a single representative scenario
(\texttt{groups\_balanced\_moderate\_drop}). All datasets are generated with
fixed seeds (split-specific offsets), so each scenario corresponds to one
deterministic dataset per split.
We use the tuning split only for model selection and hyperparameter choices;
all reported reconstruction metrics are computed on the held-out test split.

\paragraph{Scenario catalog.}
The benchmark contains 13 scenarios: one no-dropout baseline
(\texttt{groups\_balanced\_nodrop}) and six scenario families, each with an
experiment-level dropout variant and a group-level dropout variant:
balanced groups (\texttt{groups\_balanced\_moderate\_drop}),
imbalanced groups (\texttt{groups\_imbalanced\_moderate\_drop}),
rare populations (\texttt{groups\_rare\_high\_drop}),
batch effects (\texttt{batch\_effects\_moderate\_drop}),
linear trajectories (\texttt{paths\_linear\_moderate\_drop}), and
branching trajectories (\texttt{paths\_branching\_moderate\_drop}).
Group-level variants append \texttt{\_group} to the scenario id.

\paragraph{Grouped population scenarios.}
Grouped datasets use Splat's \texttt{groups} mode with group proportions
(\texttt{group.prob}): balanced $(1/3,1/3,1/3)$, imbalanced $(0.6,0.3,0.1)$, and
rare $(0.50,0.25,0.20,0.05)$. We fix the fraction of differentially expressed
genes (\texttt{de.prob}) to $0.10$ in all grouped scenarios.
Dropout is simulated using Splatter's logistic dropout model, which maps log
mean expression to a dropout probability via a logistic function with midpoint
(\texttt{dropout.mid}) and slope (\texttt{dropout.shape}); we use
moderate-dropout settings $\texttt{dropout.mid}=3.0$ and
$\texttt{dropout.shape}=-1.0$, and rare high-dropout settings
$\texttt{dropout.mid}=4.0$ and $\texttt{dropout.shape}=-1.0$.
The scope of dropout heterogeneity is controlled by \texttt{dropout.type}:
experiment-level variants use \texttt{dropout.type=experiment} (shared
parameters), while group-level variants use \texttt{dropout.type=group} with
group-specific parameter vectors created by evenly perturbing the base values
across groups (\texttt{dropout.mid} range $\pm 0.4$, \texttt{dropout.shape}
range $\pm 0.25$).

\paragraph{Batch-effect scenario.}
The batch-effect family uses two batches of approximately equal size
(\texttt{batchCells}) and introduces batch-specific multiplicative gene factors
via the batch effect location/scale parameters
(\texttt{batch.facLoc}=0.10, \texttt{batch.facScale}=0.10). We also increase
library-size dispersion (\texttt{lib.scale}=0.40) to induce greater
heterogeneity in per-cell sequencing depth.

\paragraph{Trajectory scenarios.}
Trajectory datasets are generated with Splat's \texttt{paths} mode. Both linear
and branching trajectories use four groups with equal proportions
$(0.25,0.25,0.25,0.25)$, a higher fraction of DE genes (\texttt{de.prob}=0.20),
and stronger DE factor magnitudes via the log-normal location parameter
(\texttt{de.facLoc}=0.20). We set the number of interpolation steps
(\texttt{path.nSteps}) to 50. The trajectory topology is controlled by
\texttt{path.from}: the linear topology uses $\texttt{path.from}=(0,1,2,3)$,
while the branching topology uses $\texttt{path.from}=(0,1,1,3)$. Both
trajectory families use the same moderate dropout settings as the grouped
moderate-dropout scenarios, again with experiment-level
(\texttt{dropout.type=experiment}) and group-level (\texttt{dropout.type=group})
dropout variants.

\paragraph{Assays, normalization, and marker genes.}
Each dataset stores observed counts (\texttt{counts}) and pre-dropout ground
truth (\texttt{TrueCounts}), as well as CP10k-normalized log-transformed assays
$\texttt{logcounts}=\log_2(1+\mathrm{CP10k}(\texttt{counts}))$ and
$\texttt{logTrueCounts}=\log_2(1+\mathrm{CP10k}(\texttt{TrueCounts}))$, where
observed and true counts are normalized independently. Per-cell observed and
true library sizes and the normalization target sum (\texttt{targetSum}) are
stored in metadata. Marker genes are defined using \texttt{DEFacGroup*}
row-data columns: a gene is a marker if its differential-expression factor is
not equal to 1 in any group.

